\documentclass[aspectratio=169,10pt]{beamer}
\newcommand{\reporoot}{../..}
% Shared preamble of the main deck and the topic decks.
% Define \reporoot (the path from the .tex file to the repository root)
% before \input-ing this file.
\usetheme[progressbar=frametitle,numbering=fraction,sectionpage=progressbar]{metropolis}

\usepackage{amsmath,amssymb}
\usepackage{booktabs}
\usepackage{array}
\newcolumntype{L}[1]{>{\raggedright\arraybackslash}p{#1}}
\usepackage{graphicx}
\usepackage{tikz}
\usetikzlibrary{arrows.meta,positioning,decorations.pathreplacing}
\usepackage[authoryear,round]{natbib}
\usepackage{xcolor}

\graphicspath{{\reporoot/figures/}}

% Every simulated number on the slides is produced by experiments/*.py and read from here.
\input{\reporoot/results/e01_estimators.tex}
\input{\reporoot/results/e02_scaling.tex}
\input{\reporoot/results/e03_var_backtests.tex}
\input{\reporoot/results/e04_es_backtest.tex}
\input{\reporoot/results/e05_pla.tex}
\input{\reporoot/results/e06_estimation_risk.tex}
\input{\reporoot/results/e07_model_uncertainty.tex}
\input{\reporoot/results/e08_asrf.tex}
\input{\reporoot/results/e09_copula_is.tex}
\input{\reporoot/results/e10_calibration_tests.tex}
\input{\reporoot/results/e11_discrimination.tex}
\input{\reporoot/results/e12_exposure.tex}
\input{\reporoot/results/e13_cva_wwr.tex}
\input{\reporoot/results/e14_nested.tex}
\input{\reporoot/results/e15_variance_reduction.tex}
\input{\reporoot/results/e16_lda.tex}
\input{\reporoot/results/e17_evt.tex}
\input{\reporoot/results/e18_reverse_stress.tex}
\input{\reporoot/results/e19_regimes.tex}
\input{\reporoot/results/gallery.tex}

\newcommand{\E}{\mathbb{E}}
\newcommand{\Var}{\operatorname{Var}}
\newcommand{\VaR}{\mathrm{VaR}}
\newcommand{\ES}{\mathrm{ES}}
\newcommand{\dd}{\mathrm{d}}
\definecolor{accent}{HTML}{0C5DA5}
\definecolor{warn}{HTML}{C0392B}
\newcommand{\hl}[1]{\textcolor{accent}{\textbf{#1}}}

\setbeamerfont{footnote}{size=\tiny}
\setbeamerfont{caption}{size=\scriptsize}

\author{Yaohui Shu}
\institute{Stony Brook University}
\date{November 2026}

% Topic decks: a references frame that lists only the works cited in the deck.
\newcommand{\deckreferences}{%
  \begin{frame}[allowframebreaks]{References}
    \tiny
    \renewcommand{\bibsection}{}
    \bibliographystyle{plainnat}
    \bibliography{\reporoot/slides/refs}
  \end{frame}}

% Topic decks: \decknumber{n} puts the deck number on the title page.
\newcommand{\decknumber}[1]{\institute{Stony Brook University\hfill{\scriptsize\color{gray}Deep dive #1 of 9}}}


\title{Monte Carlo Methods for Risk}
\subtitle{Nested simulation, importance sampling and quasi-Monte Carlo}
\decknumber{7}

\begin{document}

\maketitle

\begin{frame}{The question}
  \small
  \begin{columns}[T]
    \begin{column}{0.55\textwidth}
      Most risk numbers are Monte Carlo estimates, and Monte Carlo has its own error. \hl{How should a fixed computing budget be spent, and how small can the error be made?}
      \begin{itemize}
        \item Tail probabilities are rare events: plain sampling wastes most scenarios.
        \item Portfolios of derivatives need a valuation inside every scenario: simulation inside simulation.
        \item Smooth expectations converge faster with well-spread points than with random ones.
      \end{itemize}
    \end{column}
    \begin{column}{0.42\textwidth}
      \begin{alertblock}{In this deck}
        \small
        \begin{enumerate}
          \item Monte Carlo error.
          \item Nested simulation: outer versus inner samples.
          \item Importance sampling for tails.
          \item Quasi-Monte Carlo.
        \end{enumerate}
      \end{alertblock}
    \end{column}
  \end{columns}
\end{frame}

\begin{frame}{Simulation inside simulation, and better-spread points}
  \centering
  \includegraphics[width=0.95\textwidth,height=0.54\textheight,keepaspectratio]{gallery/g07_monte_carlo.png}\par\vspace{0.4em}
  \footnotesize\raggedright
  (a) Nested simulation of an underlying price $S_t$: outer scenarios up to a three-month horizon, then inner paths that reprice the position in each one. (b, c) 512 random and 512 scrambled Sobol points $(u_1, u_2)$ in the unit square, coloured by draw order. The Sobol points fill the square evenly; their L$_2$-star discrepancy is \GalSevenDiscRatio{} times smaller.
\end{frame}

\begin{frame}{Monte Carlo error, and why tails are hard}
  \small
  Estimating $\mu = \E[f(X)]$ by the mean of $n$ independent draws has standard error $s/\sqrt n$, with $s$ the standard deviation of $f(X)$ \citep{glasserman2003monte}.
  \begin{columns}[T]
    \begin{column}{0.5\textwidth}
      \textbf{Tail probability} $\alpha = P(L > c)$: indicator with variance $\alpha(1 - \alpha)$, so the relative error is
      \[ \sqrt{\frac{1 - \alpha}{\alpha\,n}} \approx \frac{1}{\sqrt{\alpha n}} . \]
      For 10\% relative error at $\alpha = 10^{-4}$ we need $10^6$ scenarios.
    \end{column}
    \begin{column}{0.46\textwidth}
      \textbf{Quantiles} inherit this: the error of an estimated VaR scales like $\sqrt{\alpha(1-\alpha)/n}$ divided by the loss density at the quantile, which is small in the tail.
      \vspace{0.3em}

      Every simulated risk number in a validation report needs its Monte Carlo error next to it.
    \end{column}
  \end{columns}
\end{frame}

\begin{frame}{Nested simulation}
  \small
  Risk at horizon $\tau$ of a derivatives portfolio needs its value $V_\tau$ in each outer scenario, which is itself an expectation. With $m$ outer scenarios and $n$ inner draws each, $\hat V_\tau = V_\tau + \text{noise}$ with variance $s^2/n$.
  \begin{columns}[T]
    \begin{column}{0.52\textwidth}
      \citet{gordy2010nested}: for $\alpha = P(L > c)$ the inner noise biases the estimate by about $\theta/n$, and the outer sampling adds variance $\alpha(1-\alpha)/m$. With budget $k = mn$,
      \[ \mathrm{MSE} \approx \frac{\theta^2}{n^2} + \frac{\alpha(1-\alpha)\,n}{k}, \]
      minimized at $n \propto k^{1/3}$, so the best RMSE falls like $k^{-1/3}$, not $k^{-1/2}$.
    \end{column}
    \begin{column}{0.45\textwidth}
      \textbf{Test portfolio.} A long one-year at-the-money call (Black-Scholes, $\sigma = 25\%$, $r = 2\%$), one-week horizon; $c = \EFourteenC$ is the exact 1\% loss level. The exact inner value is known, so the bias is measurable.
      \vspace{0.3em}

      A small inner sample overstates tail probabilities: noise widens the loss distribution.
    \end{column}
  \end{columns}
\end{frame}

\begin{frame}{Splitting the budget}
  \centering
  \includegraphics[width=0.86\textwidth]{fig14_nested.png}
  \begin{columns}[T]
    \begin{column}{0.5\textwidth}
      \footnotesize
      (a) RMSE of $\hat P(L > c)$ (true 1\%) against the inner sample size, 100 replications per point. With $n = 8$ the estimate is off by $\EFourteenBiasEight$ percentage points: inner noise swamps the one-week loss. The best $n$ rises from $\EFourteenBestNSmall$ to $\EFourteenBestNLarge$ as the budget grows 16-fold.
    \end{column}
    \begin{column}{0.46\textwidth}
      \footnotesize
      (b) Best RMSE against the budget, three budgets: fitted slope $\EFourteenSlope$. The asymptotic rate is $-1/3$; with $n$ restricted to powers of two and 100 replications, the optimum is located only roughly, so this slope is a coarse check.
    \end{column}
  \end{columns}
\end{frame}

\begin{frame}{Importance sampling and quasi-Monte Carlo}
  \small
  \begin{columns}[T]
    \begin{column}{0.5\textwidth}
      \textbf{Importance sampling.} Sample from a density $g$ that makes the event common and reweight by $f/g$. For $P(X > c)$, $X \sim \mathcal N(0,1)$, sample $X \sim \mathcal N(c, 1)$:
      \[ \hat\alpha = \frac1n\sum \mathbf{1}\{X_i > c\}\,e^{-cX_i + c^2/2}. \]
      The relative error stays bounded as $c$ grows. Portfolio credit risk uses the same idea with a factor shift and a twist (deck 4).
    \end{column}
    \begin{column}{0.46\textwidth}
      \textbf{Quasi-Monte Carlo.} Replace random points by low-discrepancy points (Sobol). Randomized (scrambled) points keep unbiasedness and give an error estimate from independent replications. For smooth integrands the error can fall close to $n^{-1}$, against $n^{-1/2}$.
      \vspace{0.3em}

      Test: arithmetic basket call on five correlated assets.
    \end{column}
  \end{columns}
\end{frame}

\begin{frame}{Variance reduction in numbers}
  \centering
  \includegraphics[width=0.86\textwidth]{fig15_variance_reduction.png}
  \begin{columns}[T]
    \begin{column}{0.5\textwidth}
      \footnotesize
      (a) $10^5$ samples, 50 replications. At $c = 4$ ($\alpha \approx 3\times10^{-5}$) plain Monte Carlo has relative error $\EFifteenRelMcFour\%$, importance sampling $\EFifteenRelIsFour\%$; at $c = 6$ importance sampling still gives $\EFifteenRelIsSix\%$, where plain Monte Carlo sees no hits at all.
    \end{column}
    \begin{column}{0.46\textwidth}
      \footnotesize
      (b) Basket price $\EFifteenRef$. Fitted error slopes: Monte Carlo $\EFifteenSlopeM$, scrambled Sobol $\EFifteenSlopeQ$. At $2^{16}$ points Sobol is $\EFifteenGain$ times more accurate.
    \end{column}
  \end{columns}
\end{frame}

\begin{frame}{Discussion and takeaways}
  \small
  \begin{columns}[T]
    \begin{column}{0.5\textwidth}
      \textbf{What the experiments show}
      \begin{itemize}
        \item In nested simulation, too few inner samples bias tail risk upward, and the optimal split follows $k^{1/3}$.
        \item Importance sampling makes rare-event estimates feasible; its quality depends on choosing the right change of measure.
        \item Quasi-Monte Carlo speeds up smooth expectations, such as prices and expected exposures, more than tail indicators.
      \end{itemize}
    \end{column}
    \begin{column}{0.46\textwidth}
      \textbf{For a risk engine}
      \begin{enumerate}
        \item Report Monte Carlo error with every risk number.
        \item For nested problems, tune the inner sample size, or replace inner simulation by a regression proxy and validate it.
        \item Use importance sampling for 99.9\% credit and operational capital.
      \end{enumerate}
      \textbf{Open question.} Error bounds for machine-learned pricing proxies inside risk engines.
    \end{column}
  \end{columns}
\end{frame}

\begin{frame}[fragile]{Reproduce}
  \small
  Code: \texttt{montecarlo.py}. Experiments: \texttt{e14\_nested.py}, \texttt{e15\_variance\_reduction.py}.
\begin{verbatim}
from stochastic_risk import montecarlo as mc
c = mc.loss_threshold(0.01, S0=100, K=100, r=0.02, sigma=0.25, T=1.0, tau=1/52)
mc.nested_tail_probability(m=4096, n=64, c=c, S0=100, K=100, r=0.02,
                           sigma=0.25, T=1.0, tau=1/52)
mc.gaussian_tail_is(c=4.0, n=100000)
\end{verbatim}
\end{frame}

\deckreferences
\end{document}
